% Part: turing-machines 
% Chapter: machines-computations 
% Section: introduction

\documentclass[../../../include/open-logic-section]{subfiles}

\begin{document}

\olfileid[hi]{tur}{mac}{int}
\olsection{परिचय}

मान लें कि कोई फलन $\Nat$ से $\Nat$ में जाता है। उसके
\emph{संगणनीय} होने का क्या अर्थ है? सबसे आरंभिक और सर्वाधिक प्रसिद्ध
उत्तरों में से एक यह है कि कोई फलन तभी संगणनीय है जब उसे किसी ट्यूरिंग
मशीन द्वारा संगणित किया जा सके। एलन ट्यूरिंग ने यह धारणा 1936 में
प्रस्तुत की थी। ट्यूरिंग मशीनें \emph{संगणना के एक प्रतिरूप} का उदाहरण
हैं---वे ``संगणनात्मक प्रक्रिया'' के विचार को परिभाषित करने की एक
गणितीयतः सटीक रीति देती हैं। इसका ठीक-ठीक अर्थ क्या है, इस पर मतभेद है;
फिर भी इस बात पर व्यापक सहमति है कि ट्यूरिंग मशीनें संगणनात्मक
प्रक्रियाओं को निर्दिष्ट करने की एक रीति हैं। यद्यपि ``ट्यूरिंग मशीन''
पद चलायमान पुर्जों वाली किसी भौतिक मशीन की छवि उत्पन्न करता है, किंतु
कड़ाई से कहें तो ट्यूरिंग मशीन पूर्णतः गणितीय रचना है और इस रूप में वह
संगणनात्मक प्रक्रिया के विचार का आदर्शीकरण करती है। उदाहरणार्थ, हम
ट्यूरिंग मशीन की समय अथवा स्मृति-संबंधी आवश्यकताओं पर कोई प्रतिबंध नहीं
लगाते: ट्यूरिंग मशीन किसी वस्तु को तब भी संगणित कर सकती है जब उस संगणना
के लिए ब्रह्मांड में उपस्थित परमाणुओं की संख्या से अधिक भंडारण-स्थान या
अधिक चरणों की आवश्यकता पड़े।

\begin{explain}
ट्यूरिंग मशीन को किसी विशेष प्रकार की काल्पनिक युक्ति के प्रोग्राम के
रूप में समझना संभवतः सर्वोत्तम है। इस युक्ति में एक \emph{पट्टी} और एक
\emph{पठन-लेखन शीर्ष} होता है। ट्यूरिंग मशीन के हमारे संस्करण में पट्टी
एक दिशा में (दाईं ओर) अनंत होती है और \emph{खानों} में विभाजित होती है;
प्रत्येक खाने में किसी परिमित \emph{वर्णमाला} का एक प्रतीक हो सकता है।
ऐसी वर्णमालाओं में कितने भी भिन्न प्रतीक हो सकते हैं, किंतु हम मुख्यतः
तीन प्रतीकों से काम चलाएँगे: $\TMendtape$, $\TMblank$, और $\TMstroke$।
युक्ति को आरंभ करते समय पट्टी रिक्त होती है (अर्थात् प्रत्येक खाने में
$\TMblank$ प्रतीक होता है), सिवाय सबसे बाएँ खाने के, जिसमें
$\TMendtape$ होता है, और उन परिमित खानों के जिनमें \emph{आगत} होता है।
किसी भी समय युक्ति परिमित संख्या में उपलब्ध \emph{अवस्थाओं} में से एक
में होती है। आरंभ में शीर्ष सबसे बाएँ खाने को पढ़ता है और युक्ति एक
निर्दिष्ट \emph{आरंभिक अवस्था} में होती है। युक्ति की क्रिया के प्रत्येक
चरण में वर्तमानतः पढ़े जा रहे खाने की सामग्री, युक्ति की अवस्था और
ट्यूरिंग मशीन का प्रोग्राम मिलकर निर्धारित करते हैं कि आगे क्या होगा।
ट्यूरिंग मशीन का प्रोग्राम एक आंशिक फलन द्वारा दिया जाता है, जो किसी
अवस्था~$q$ और प्रतीक~$\sigma$ को आगत के रूप में लेकर त्रिक
~$\tuple{q', \sigma',
  D}$ निर्गत करता है। जब भी युक्ति अवस्था $q$ में
होती है और प्रतीक $\sigma$ पढ़ती है, वह वर्तमान खाने के प्रतीक को
$\sigma'$ से बदल देती है; $D$ क्रमशः $\TMleft$, $\TMright$, या
$\TMstay$ होने के अनुसार शीर्ष बाएँ जाता है, दाएँ जाता है, या वहीं रहता
है; और युक्ति अवस्था~$q'$ में चली जाती है।

उदाहरणार्थ, \olref{fig:tm} की स्थिति पर विचार करें।
\begin{figure}
  \olasset{\olpath/assets/diagrams/turing-machine.tikz}
  \caption{अपना प्रोग्राम निष्पादित करती हुई एक ट्यूरिंग मशीन।}
  \ollabel{fig:tm}
\end{figure}
ट्यूरिंग मशीन की पट्टी के दिखाई देने वाले भाग में सबसे बाएँ खाने पर
पट्टी-अंत प्रतीक $\TMendtape$ है; उसके बाद तीन $1$, एक $0$, और फिर चार
$1$ हैं। शीर्ष बाएँ से तीसरे खाने को पढ़ रहा है, जिसमें एक~$1$ है, और
वह अवस्था~$q_1$ में है---हम कहते हैं, ``मशीन $1$ को अवस्था~$q_1$ में पढ़
रही है।'' यदि ट्यूरिंग मशीन का प्रोग्राम आगत $\tuple{q_1, 1}$ के लिए
त्रिक $\tuple{q_2,
0, \TMstay}$ लौटाता है, तो मशीन तीसरे खाने के~$1$ को
$0$ से बदल देगी, पठन-लेखन शीर्ष को वहीं रहने देगी और अवस्था~$q_2$ में
चली जाएगी। यदि इसके बाद प्रोग्राम $\tuple{q_3, 0, \TMright}$ को आगत
$\tuple{q_2, 0}$ के लिए लौटाता है, तो मशीन उस~$0$ के ऊपर एक और~$0$
लिखेगी (अर्थात् पठन-लेखन शीर्ष के नीचे पट्टी की सामग्री वस्तुतः अपरिवर्तित
रहेगी), एक खाना दाईं ओर जाएगी और अवस्था~$q_3$ में प्रवेश करेगी। इसी प्रकार
क्रिया आगे चलती रहती है।

हम कहते हैं कि मशीन तब \emph{रुकती} है जब वह किसी ऐसी अवस्था $q_n$ और
प्रतीक $\sigma$ पर पहुँचती है जिसके लिए $\tuple{q_n, \sigma}$ का कोई
निर्देश नहीं है; अर्थात् आगत $\tuple{q_n,\sigma}$ पर संक्रमण फलन
अपरिभाषित है। दूसरे शब्दों में, मशीन के पास निष्पादित करने के लिए कोई
निर्देश नहीं होता और उस बिंदु पर उसकी क्रिया समाप्त हो जाती है। रुकने को
कभी-कभी एक विशिष्ट विराम-अवस्था~$h$ द्वारा दर्शाया जाता है। इसे आगे अधिक
विस्तार से दिखाया जाएगा।
\end{explain}

\begin{digress}
ट्यूरिंग के शोधपत्र ``संगणनीय संख्याओं पर'' की सुंदरता यह है कि उसमें
वे केवल एक आकारिक परिभाषा ही नहीं देते, बल्कि यह तर्क भी देते हैं कि वह
परिभाषा संगणनीयता की सहज धारणा को ग्रहण करती है। परिभाषा से यह स्पष्ट
होना चाहिए कि ट्यूरिंग मशीन द्वारा संगणनीय कोई भी फलन सहज अर्थ में
संगणनीय है। ट्यूरिंग इस विलोम के सत्य होने के पक्ष में तीन प्रकार के
तर्क देते हैं; अर्थात् जिस किसी फलन को हम स्वाभाविक रूप से संगणनीय
मानेंगे, वह ऐसी मशीन द्वारा संगणनीय है। वे (ट्यूरिंग के शब्दों में) हैं:
\begin{enumerate}
\item सहज बोध की सीधी दुहाई।
\item दो परिभाषाओं की समतुल्यता का प्रमाण (यदि नई परिभाषा सहज बोध को
  अधिक आकर्षित करती हो)।
\item संगणनीय संख्याओं के बड़े वर्गों के उदाहरण देना।
\end{enumerate}
हमारा लक्ष्य संगणनीयता की धारणा को ``सिद्धांततः'' परिभाषित करना है;
अर्थात् समय और स्थान की व्यावहारिक सीमाओं को ध्यान में रखे बिना। निश्चय
ही, संगणनीयता की व्यापकतम परिभाषा स्थापित हो जाने पर परिबद्ध संसाधनों
वाली संगणना पर विचार किया जा सकता है; यही ``संगणनात्मक जटिलता'' नामक
विषय का केंद्र है।
\end{digress}

\begin{history}
एलन ट्यूरिंग ने 1936 में ट्यूरिंग मशीनों का आविष्कार किया। उस समय उनकी
रुचि प्रथम-कोटि तर्क की निर्णेयता में थी, फिर भी इस शोधपत्र को संगणक की
अभिकल्पना की नींव पर एक निर्णायक शोधपत्र कहा गया है। उसमें ट्यूरिंग का
ध्यान संगणनीय वास्तविक संख्याओं पर है, अर्थात् ऐसी वास्तविक संख्याएँ
जिनके दशमलव प्रसार संगणनीय हैं; किंतु वे बताते हैं कि उनकी धारणाओं को
प्राकृतिक संख्याओं पर संगणनीय फलनों आदि के अनुकूल बनाना कठिन नहीं है।
ध्यान दें कि यह 1941 में पहले कार्यशील सामान्य-प्रयोजन संगणक के निर्माण
से पूरे पाँच वर्ष पहले था (जर्मन कोनराड त्सूज़े ने इसे अपने माता-पिता के
बैठक-कक्ष में बनाया था); ब्लेच्ली पार्क में ट्यूरिंग और उनके सहकर्मियों
द्वारा कूट-भेदी कोलोसस (1943) बनाए जाने से सात वर्ष पहले; अमेरिकी
ईएनिऐक (1945) से नौ वर्ष पहले; मैनचेस्टर में पहले ब्रिटिश सामान्य-प्रयोजन
संगणक---मैनचेस्टर लघु-माप प्रायोगिक मशीन---के निर्माण (1948) से बारह
वर्ष पहले; और अमेरिकियों द्वारा बाइनैक (1949) के प्रथम परीक्षण से तेरह
वर्ष पहले। मैनचेस्टर एसएसईएम को प्रथम संग्रहित-प्रोग्राम संगणक होने का
गौरव प्राप्त है---इससे पहले की मशीनों को प्रत्येक नए कार्य के लिए हाथ से
पुनः तारबद्ध करना पड़ता था।
\end{history}

\end{document}
